% $Id: INF_Poster_example.tex 7714 2011-08-31 17:34:46Z tkren $
%
% TU Wien - Faculty of Informatics
% poster template
%
% This template is using the beamer document class and beamerposter package, see
% <http://www.ctan.org/tex-archive/macros/latex/contrib/beamer/>
% <http://www.ctan.org/tex-archive/macros/latex/contrib/beamerposter/>
% <http://www-i6.informatik.rwth-aachen.de/~dreuw/latexbeamerposter.php>
%
% For questions and comments send an email to
% Thomas Krennwallner <tkren@kr.tuwien.ac.at>
%

\documentclass[final,hyperref={pdfpagelabels=true}]{beamer}

\usepackage{TUINFPST}

\usepackage{lipsum}
 
%\title[Computational Intelligence]{Interactive Computer Generated Architecture}
% if you have a long title looking squeezed on the poster, just force
% some distance:
\title[Computational Intelligence]{%
  Interpolation in First-Order Logic
	\\[0.2\baselineskip]%
	with Equality
  %\\[0.2\baselineskip]%
}
%\author[b.mallinger@gmx.at]{Bernhard Mallinger}
\author[b.mallinger@gmx.at]{Bernhard Mallinger}
\institute[]{%
  Technische Universit{\"a}t Wien\\[0.25\baselineskip]
  Institut f{\"u}r diskrete Mathematik und Geometrie\\[0.25\baselineskip]
  Arbeitsbereich: Computational Logic\\[0.25\baselineskip]
  Betreuer: Ass.Prof.~Stefan Hetzl
}
%\titlegraphic{ {\Huge \color{red} TODO:~LOGO}\includegraphics[height=52mm]{logo_KBS_2_CMYK}}
\date[\today]{\today}
\subject{epilog}
\keywords{my kwd1, my kwd2}

%%%%%%%%%%%%%%%%%%%%%%%%%%%%%%%%%%%%%%%%%%%%%%%%%%%%%%%%%%%%%%%%%%%%%%%%%%%%%%%%%%%%%%

% Display a grid to help align images 
%\beamertemplategridbackground[12.7mm]

% play around with the background colors
% \setbeamercolor{background canvas}{bg=yellow}

% use a background picture
% \usebackgroundtemplate{%
%   \includegraphics[width=\paperwidth]{logo_KBS_2_CMYK}
% }

% play around with block colors
\setbeamercolor{block body}{fg=black,bg=white}
\setbeamercolor{block title}{fg=TuWienBlue,bg=white}

\setbeamertemplate{block begin}{
  \begin{beamercolorbox}{block title}%
    \begin{tikzpicture}%
      \node[draw,rectangle,line width=3pt,rounded corners=0pt,inner sep=0pt]{%
        \begin{minipage}[c][2cm]{\linewidth}
          \centering\textbf{\insertblocktitle}
        \end{minipage}
      };
    \end{tikzpicture}%
  \end{beamercolorbox}
  \vspace*{1cm}
  \begin{beamercolorbox}{block body}%
}

\setbeamertemplate{block end}{
  \end{beamercolorbox}
  \vspace{2cm}
}

% setup postit
\setbeamercolor{postit}{fg=black,bg=yellow} 
\newenvironment{postit}
{\begin{beamercolorbox}[sep=1em,wd=7cm]{postit}}
{\end{beamercolorbox}}

\usepackage{tikz}
\usetikzlibrary{shapes,arrows,backgrounds,graphs,%
	matrix,patterns,arrows,decorations.pathmorphing,decorations.pathreplacing,%
	positioning,fit,calc,decorations.text,shadows,decorations.markings%
}


% for crop marks, uncomment the following line
%\usepackage[cross,width=88truecm,height=123truecm,center]{crop}

%%%%%%%%%%%%%%%%%%%%%%%%%%%%%%%%%%%%%%%%%%%%%%%%%%%%%%%%%%%%%%%%%%%%%%%%%%%%%%%%%%%%%%

\begin{document}

% We have a single poster frame.
\begin{frame}
	\begin{columns}[t]


		\begin{column}{.45\textwidth}

			\begin{block}{Interpolation}
				\begin{itemize}

					\item
						The notion of interpolation has been introduced by Craig.

					\item
						Loosely speaking, given two formulas $A$ and $B$ such tht $A$ implies $B$, an interpolant $I$ is a formula which is implied by $A$ and which itself implies $B$.
						\bigskip


						\begin{figure}[htbp]
							\centering
							\begin{tikzpicture}[
									implies/.style={double,double equal sign distance,-implies},
									mynode/.style={circle,outer sep=3pt}
								]
								\node[mynode] (A) at (0,0) {$A$};
								\node[mynode] (B) at (16,0) {$B$};
								\node[mynode] (I) at (8,-6) {$I$};

								%\draw[->,implies] (A) to (B);
								\draw (A) edge[implies]  (B);
								\draw (A) edge[implies]  (I);
								\draw (I) edge[implies]  (B);

							\end{tikzpicture}
							\label{fig:interpol}
							%\caption{Given two formulas $A$ and $B$ such that $A$ implies $B$, an interpolant is a formula $I$ which is implied by $A$ and which implies $B$.}
						\end{figure}
						\smallskip

					\item
						Hence it in some sense captures the logical content of $A$ which necessarily makes $B$ true and therefore acts as a link between these formulas.

					\item
						Additionally, the language of an interpolant must be contained in the intersection of the langauges $A$ and $B$.

				\end{itemize}

			\end{block}


			\begin{block}{Aim of thesis}
				Give comprehensive account of existing proofs:
				\begin{itemize}
					\item foo
				\end{itemize}
			\end{block}


		\end{column}

		% ---------------------------------------------------------%
		% Set up a column
		%\begin{column}{.45\textwidth}
		%  \begin{block}{Block 1}
		%    \lipsum[1-3]
		%  \end{block}

		%  \begin{block}{Block 2}
		%    \lipsum[3-4]
		%  \end{block}
		%\end{column}
		% ---------------------------------------------------------%
		% end the column

		% ---------------------------------------------------------%
		% Set up a column 
		\begin{column}{.45\textwidth}
			\begin{block}{Yet another block}
				\lipsum[5-7]
			\end{block}



			\begin{block}{References}
				% this is just an example, use BibTeX!
				\begin{thebibliography}{999}

					\bibitem[1]{Craig57linear}
						William Craig. Linear Reasoning. A New Form of the Herbrand-Gentzen Theorem.
						\newblock {
							\emph{Journal of Symbolic Logic}, 22(3):250--268, 1957.
						}

					\bibitem[2]{Huang95}
						Guoxiang Huang. Constructing Craig Interpolation Formulas.
						\newblock {
							\emph{Proceedings of the First Annual International Conference on Computing and Combinatorics}, COCOON ’95, p.\ 181--190, 1995. 
						}

						%  \newblock {On logical representations of hackerisms}.
						%  {\em J.~Log.~Hack.} 1:1--2.

						%	
						%  \bibitem[Foo~and~Fu, 2010]{ff2010}
						%    Foo, B.; and Fu, B.
						%    2010.
						%    \newblock {On logical representations of hackerisms}.
						%    {\em J.~Log.~Hack.} 1:1--2.
						%    
						%    
						%  \bibitem[Crock~et~al., 2010]{ck2010}
						%    Crock, A; Cruft, B.; and Kludge, C.
						%    2010.
						%    \newblock {Decomposing junk code}.
						%    Manuscript.
						%    
				\end{thebibliography}
			\end{block}
		\end{column}
		% ---------------------------------------------------------%
		% end the column
	\end{columns}

	%\begin{tikzpicture}[remember picture,overlay]
	%	\node[inner sep=0pt,xshift=-30cm,yshift=23cm] at (current page.east) {%
	%		\begin{postit}%
	%			Post-It time!%
	%		\end{postit}%
	%	}; 
	%\end{tikzpicture}

\end{frame}

\end{document}

%%% Local Variables:
%%% TeX-PDF-mode: t
%%% TeX-debug-bad-boxes: t
%%% TeX-master: t
%%% TeX-parse-self: t
%%% TeX-auto-save: t
%%% reftex-plug-into-AUCTeX: t
%%% End:
